\section*{The Propers: Detailed Commentary}

\subsection*{Tobit's Mercy Refrain: Received, Returned, and Proclaimed\enspace\properrefs{Int., Off., Comm.}}

The missal cites Tobit 12:6 three times, but it does not merely repeat one antiphon. Raphael's exhortation follows the family's restored sight, recovered household, and offer of generous payment. He redirects father and son from the angelic servant to the God whose will governed the journey: bless God, give thanks before all the living, proclaim his deeds, and write what happened. Mercy is therefore remembered as a gift that asks for truthful public speech.

The Roman chants give that command an explicitly Trinitarian form, which is the liturgy's adaptation and not the wording of Tobit itself.

\begin{threestable}{Liturgical moment}{What changes in the chant}{Resulting action}
Introit & Holy Trinity and undivided Unity are blessed; the assembly promises, \latin{confitebimur ei}, because mercy has been shown. & Received mercy gives the entering Church her first-person plural voice before doctrine is analyzed.\\
Offertory & Father, only-begotten Son, and Holy Ghost are named distinctly while the same mercy clause remains. & Praise becomes more explicit as the gifts approach the altar; confession accompanies offering rather than replacing it.\\
Communion & The title ``God of heaven'' returns, blessing is present-tense, and confession is placed \latin{coram omnibus viventibus}. & Sacramental reception opens toward testimony before all the living instead of ending in private possession.\\
\end{threestable}

St.~Thomas Aquinas receives the verse itself, exactly though briefly. In his \emph{Expositio in Psalmos} on Psalm 40, no.~9, he cites \latin{Tob.~12: benedicite Deum caeli} while distinguishing the way creatures bless God---by recognizing his goodness and confessing it orally---from the way God blesses creatures by causing goodness in them. Bede's \emph{Allegorica interpretatio in Tobiam} is a second medieval witness to the chapter, though not to the appointed imperative itself. At Tobit XII (PL 91:935B--C), Bede receives Raphael's disclosure and return as a figure of fuller divine self-disclosure: Christ is equal to the Father in divinity and consubstantial with humanity. Aquinas's exact citation and Bede's chapter-level reading illuminate why revelation issues in praise, though neither supplies a continuous commentary on Tobit 12:6.

Tobit's surrounding exhortation joins thanksgiving to prayer, fasting, almsgiving, burial of the dead, and the family's concrete healing. The proper draws none of these details into an allegory of this Mass. Its direct force is simpler and demanding: divine beneficence precedes the response, and praise of that beneficence must become visible in truthful speech and a life that honors the gift.

\begin{studybox}{Two forms of Tobit 12}
The Vulgate form cited by the missal and a modern critical translation divide and phrase Tobit 12 differently. The appointed Latin is the 1962 proper's, and the whole chapter is its context; the chants' explicit naming of Father, Son, and Holy Ghost belongs to the liturgy, not to the earlier compositional horizon of Tobit.
\end{studybox}

\sectionguard
\subsection*{The Collect: Confession Ordered to Adoration\enspace\properrefs{Coll.}}

The Collect supplies a compact rule for approaching the feast. Its grammar begins with what God has given his servants, distinguishes acknowledgment from adoration, and asks that the same faith become a defense. Intellectual confession is neither self-generated mastery nor an end in itself.

\begin{threestable}{Latin hinge}{Work assigned to faith}{Precise sense}
\latin{qui dedisti famulis tuis} & True faith and the capacity to confess begin as divine gift to servants. & The worshipper does not manufacture the mystery or stand over it as judge.\\
\latin{gloriam agnoscere} & Faith truly acknowledges the glory of the eternal Trinity. & Incomprehensibility does not make revealed confession empty or merely emotional.\\
\latin{adorare unitatem} & Knowledge passes into worship of the Unity in the power of majesty. & Unity is adored without erasing the three Persons named by the Gospel and Preface.\\
\latin{eiusdem fidei firmitate} & The prayer asks that steadfastness in this same faith guard against every adversity. & Protection is requested through persevering faith, not promised as immunity from suffering or inquiry.\\
\end{threestable}

The Collect thus anticipates the entire order that follows. Romans permits sustained inquiry but ends in doxology; Matthew gives the baptismal Name and commands observance; the Postcommunion asks that confession continue to profit body and soul. The prayer's firmness is not rigidity against evidence. It is constancy within the truth received, worshipped, taught, and lived.

\sectionguard
\subsection*{Romans 11: Doxology after the Mercy Argument\enspace\properrefs{Ep.}}

The Epistle is only four verses, but its argument is not contextless. Paul has spent three chapters grieving over Israel, refusing Gentile boasting, distinguishing unbelief from irrevocable divine gifts, and tracing disobedience toward the declaration that God has enclosed all in disobedience so as to have mercy upon all. Only then does he cry, ``O the depth.'' The doxology marks the limit of an argument actually pursued, not a prohibition on theological reasoning.

\begin{fourstable}{Romans movement}{Question faced}{Precise sense}{Doxological result}
11:25--27 & Israel's partial hardening, the nations' fullness, and God's covenantal promise. & Gentile inclusion remains within God's covenantal fidelity to Israel. & Divine fidelity exceeds the horizon of either community's boasting.\\
11:28--32 & Enmity and election, disobedience and mercy. & Paul preserves both moral responsibility and the priority of mercy. & The saving economy ends in wonder at divine wisdom and mercy.\\
11:33--35 & Judgments, ways, counsel, and prior gift. & The creature neither maps every divine disposition nor places God in debt. & Chrysostom directly reads Paul as knowing that there is a depth without claiming its measure.\\
11:36 & All things are from, through, and unto God. & The verse remains doxological in its immediate grammar, which does not divide it among the Persons. & Inquiry resolves in glory offered to the God whose gifts precede every return.\\
\end{fourstable}

Chrysostom's direct exposition sharpens both sides. Paul has perceived a real portion of the dispensation and can proclaim its goodness; he has not searched it to exhaustion. God gives as the first benefactor, not as one repaying a prior creditor. Applied to the formulary, that reading binds the Epistle to Tobit's mercy refrain and the Secret's oblation, though Chrysostom did not comment on their later liturgical conjunction.

Aquinas's medieval lecture preserves Paul's argumentative order with unusual precision. Romans 11:33 is wonder after an attempted account of mercy toward Jews and Gentiles; ``unsearchable'' means that creatures cannot comprehend God's judgments and modes of action, not that revealed effects are unknowable. At verses 35--36, Aquinas distinguishes the giver who owes no prior return from every creature whose gift is already received from God, then reads \latin{ex}, \latin{per}, and \latin{in} as a threefold account of divine causality. This expands Chrysostom's benefactor theme without mechanically assigning the prepositions to three Persons.

\begin{studybox}{Wonder as disciplined inquiry}
``Unsearchable'' does not mean that Scripture, doctrine, and salvation history disclose nothing. Paul's immediate argument explicitly warns the grafted branches against boasting. Wonder is disciplined knowledge before divine freedom and fidelity.
\end{studybox}

\sectionguard
\subsection*{Psalm 8 and Daniel: Creaturely Praise before the Gospel\enspace\properrefs{Int., Grad., All.}}

Psalm 8 and Daniel 3 teach praise from different human positions. The Introit's psalm verse looks across the whole earth and blesses the admirable divine Name. Augustine reads that verse through Christ's humiliation and exaltation, then explains how praise from those still nourished in faith belongs to the Church's growth toward mature understanding. Beginning with praise is therefore not a refusal to learn; it is the creature's truthful entrance into learning.

Daniel's song begins inside imperial violence. Before the appointed verses, Azariah confesses sin, the absence of temple sacrifice, and the offering of a contrite heart; the three then sing with one voice in the furnace. Their blessing moves from the God of the fathers and his holy Name to temple, throne, depths, heaven, and finally the whole created order.

The Song's ancient reception survives around the appointed verses. Hippolytus's catena fragment at verse 87 marvels that the youths enumerate the whole created order and thereby confess every creature as servant of the God who created all through the Word. Jerome directly comments on verses 57--58 as the summons of creation and carefully records the Prayer and Song's absence from the Hebrew/Aramaic textual form and reception through Theodotion. Hippolytus and Jerome receive the Song's cosmic praise and its textual history; neither expounds the appointed verses 52 or 55--56, and none of the ancient witnesses checked does.

\begin{threestable}{Appointed chant}{Canonical horizon}{Preparation for what follows}
Introit, Psalm 8:2 & The Lord's Name is admirable throughout the earth, while human littleness receives a vocation under divine majesty. & Earth-wide praise prepares for Christ's command concerning all nations without turning human dignity into autonomy.\\
Gradual, Daniel 3:55--56 & God sees the depths, reigns above the cherubim, and is blessed in the firmament while the singers remain inside the furnace. & Paul's unsearchable depth is placed beneath divine sight; trial is real but never outside God's rule.\\
Alleluia, Daniel 3:52 & The God of the fathers and his glorious Name are blessed before the song summons all creation. & Universal mission carries Israel's received praise to the nations.\\
\end{threestable}

The texts establish a shared posture: the earth, the abyss, the furnace, and the heavens are all creaturely locations from which God may be blessed. The Gospel then supplies what the chants do not generate for themselves---the singular baptismal Name of Father, Son, and Holy Ghost and the command to teach obedience under Christ's continuing presence.

\sectionguard
\subsection*{Matthew 28: One Name, Three Inseparable Persons\enspace\properrefs{Gosp.}}

The Gospel's grammar holds unity and distinction together. The singular ``Name'' names Father, Son, and Holy Ghost; baptism is followed by teaching observance; the command to all nations rests on the risen Son's universal authority and continuing presence. The Fathers use this text not as an isolated proof but as a rule for Christian initiation, confession, worship, and life.

\begin{fourstable}{Witness and locus}{Direct reading of Matthew 28}{Doctrinal contribution}{Context and emphasis}
St. Hilary, \emph{De Trinitate} II.1--2 & Hilary calls the command a comprehensive saving rule whose words, actions, and sequence join the Creator, the Only-begotten, and the Gift. & The three are not separable stages or unequal powers; baptism introduces the believer into their supreme union. & Hilary immediately admits the poverty of theological language: heresy makes speech necessary, but speech never contains the mystery.\\
St. Basil, \emph{De Spiritu Sancto} 10.24, 26; 12.28; 27.68 & The Spirit's coordination with Father and Son in baptism cannot be reconciled with demoting him to a creaturely rank. Faith and baptism are completed in the same named confession. & The baptismal formula governs equal confession and doxology; saving regeneration is from Father, Son, and Holy Ghost without dividing the Spirit from their communion. & Basil's argument answers anti-Spirit controversy. It is more than a bare observation that three nouns occur in one sentence.\\
St. Gregory Nazianzen, \emph{Oration 40}.41, 45 & Gregory confesses one Godhead and power in three, then returns to Christ's command to disciple and baptize in Father, Son, and Holy Ghost. & The three share the common Name of Godhead; works must be built upon this dogmatic foundation. & Gregory does not oppose doctrine to conduct. Baptismal confession requires a life fitted to the gift received.\\
\end{fourstable}

\begin{massbox}{The appointed Preface as liturgical confession}
The Preface of the Most Holy Trinity confesses one Lord not in the singularity of one Person but in Trinity of one substance: distinction belongs to the Persons, unity to essence, equality to majesty. This is the feast's appointed liturgical context, not another row of the proper; Hilary, Basil, and Gregory expound Matthew 28, not the formulary.
\end{massbox}

The order ``make disciples, baptizing, teaching them to observe'' prevents two opposite reductions. Baptism is not a badge detached from formation, and moral formation is not a program detached from incorporation into the revealed Name. Gregory's call to build good works upon dogma closely matches the Gospel's own movement from initiation to observance. Christ's final promise keeps both mission and obedience within presence: the sender remains with the sent until the consummation.

Aquinas's \emph{Catena Aurea} shows how a medieval Doctor received this passage by disciplined collation. Its named excerpts retain Jerome's order---teach, wash in faith, then teach observance---Hilary's insistence that salvation is fully contained in the baptismal confession, Bede's reading of the mountain as resurrection height, and Chrysologus's new creation of the nations. In the catena the witnesses remain identifiable and sometimes distinct; it is Aquinas's compilation, not a fifth Father, and its excerpts stand as he transmits them.

\sectionguard
\subsection*{The Secret: The Name, the Victim, and the Offerers\enspace\properrefs{Sec.}}

The Secret is not a generic prayer over gifts. Its syntax makes a compact sacramental argument: the holy Name is invoked over the victim, and through it the worshippers ask to be perfected as God's eternal gift. The feminine pronoun \latin{eam} may look back to the victim or the oblation; the prayer's movement holds on either reading of that narrower antecedent.

\begin{threestable}{Latin hinge}{Immediate force}{Theological movement}
\latin{per tui sancti nominis invocationem} & Sanctification is asked through invocation of God's holy Name. & The Name confessed in baptism is invoked at the altar; neither use is a magical sound detached from faith and divine action.\\
\latin{huius oblationis hostiam} & The direct object of \latin{sanctifica} is the victim of this oblation. & The prayer first asks God to sanctify what is offered, not to certify an autonomous human achievement.\\
\latin{et per eam nosmetipsos} & ``And through it, us ourselves'': the petition passes through the oblation to the offerers. & External gift and interior self-offering are joined without being confused.\\
\latin{tibi perfice munus aeternum} & God is asked to perfect the worshippers for himself as an eternal gift. & The altar aims beyond a momentary sentiment toward an ecclesial, enduring consecration.\\
\end{threestable}

\begin{sourcecard}{St. Augustine --- \emph{City of God} X.5--6, 20}
Augustine is not commenting on this Secret, but his doctrine of sacrifice bears precisely on it: visible sacrifice signifies the invisible offering that unites a person to God; mercy is true sacrifice when ordered to that union; body and soul are offered as the mind is transformed; and the whole redeemed city is offered through Christ the High Priest. In the altar's sacrament the Church learns to offer herself through the one who is both priest and sacrifice.
\end{sourcecard}

Gregory Nazianzen's baptismal exhortation in \emph{Oration 40}.25 provides an analogue from Christian initiation: the baptized is to offer himself and clothe himself with Christ, making conduct answer the sacramental gift. Gregory does not comment on the Secret, but his exhortation helps explain why invocation of the baptismal Name can lead to a petition concerning \latin{nosmetipsos}.

\begin{studybox}{Sacrifice without collapse}
The victim and the congregation are not interchangeable objects, and the Church does not become Christ by nature. Christ's unique priestly self-offering makes ecclesial self-offering possible. Augustine's account preserves the order: the many become one body in the High Priest and, through him, offer themselves to God.
\end{studybox}

\sectionguard
\subsection*{Communion and Postcommunion: Confession for Body and Soul\enspace\properrefs{Comm., Postcomm.}}

The Communion returns from the altar to Raphael's command: bless the God of heaven and confess his mercy before all the living. Because the antiphon is sung at sacramental reception, public praise is not an optional social consequence attached to an otherwise private encounter. The received gift forms a confessing people. Yet the chant does not equate publicity with evangelization by technique; witness remains thanksgiving directed first to God.

The Postcommunion then joins two realities in one petition: reception of this sacrament and confession of the eternal Trinity and indivisible Unity. Both are asked to \latin{proficiat nobis}, to advance or profit the worshippers, \latin{ad salutem corporis et animae}. The prayer neither treats confession as a verbal appendage to sacrament nor treats sacramental reception as self-executing apart from faith and continuing transformation.

\begin{threestable}{Petitioned reality}{What the prayer affirms}{Reduction excluded}
\latin{huius sacramenti susceptio} & A real sacramental reception is ordered toward the salvation of the receiver. & Reception is not reduced to memory, sentiment, or a purely communal signal.\\
\latin{Trinitatis ... Unitatis confessio} & The faith confessed at baptism and throughout the feast remains integral after Communion. & Doctrine is not a decorative formula detached from sacramental and moral life.\\
\latin{salutem corporis et animae} & The petition embraces the human person bodily and spiritually. & The body is neither ignored as religiously irrelevant nor promised exemption from mortality and suffering.\\
\end{threestable}

This final oration also keeps the whole sequence from closing too quickly. The Mass does not say that one celebration has completed the offerers' perfection. It asks for progress. The feast's movement from mercy to confession, offering, reception, and witness is therefore both real and unfinished, dependent at every stage upon the God first blessed at the entrance.

\sectionguard
\subsection*{Charity: Participation, Not a Diagram\enspace\properrefs{Gosp., Sec., Comm., Postcomm.}}

Augustine's \emph{De Trinitate} VIII.7.10--VIII.10.14 moves from righteous love of the brother to the confession that God is love, and then to lover, beloved, and love. Read in full, the passage is more disciplined than its memorable phrases sometimes suggest.

\begin{threestable}{Augustine's step}{Spiritual claim}{Precise sense}
Love of the righteous brother & True love includes readiness to give oneself for the brethren and makes the unseen God known through the charity he gives. & Sentiment without righteousness or self-gift is not Augustine's charity.\\
God is love & Brotherly love is from God and gives real participation in divine life. & Created charity is not the divine essence captured as a psychological object.\\
Lover, beloved, love & Augustine identifies a genuine triadic starting point for contemplation. & He explicitly treats it as a beginning of inquiry, not a discovery that contains or proves the Trinity.\\
\end{threestable}

The Mass therefore refuses to leave Trinity at the level of an elegant schema. Baptism obliges observance; the Secret asks that the offerers themselves become gift; Communion returns to public thanksgiving; and the Postcommunion asks that sacramental reception and confession advance the salvation of body and soul. Charity is the graced form in which confessed truth becomes lived communion, not an alternative to the singular Name and three Persons revealed by Christ.

\Needspace{18\baselineskip}
\section*{The Propers: Notable and Quotable}

{\small
\begin{itemize}
\item \textbf{``O the depth'' --- John Updike, ``Tsokadze---O Altitudo,'' \emph{The New Yorker}, 25 Feb.~1956, p.~115.} The Latin incipit of Paul's astonishment becomes a learned altitude pun on a Soviet ski-jumper: ascent, descent, and thinning air turn theological depth into athletic height. The allusion is highly legible, but the publisher does not document Updike's dependence on Romans.

\item \textbf{``Who hath known the mind of the Lord?'' --- H.~L. Mencken, \emph{Damn! A Book of Calumny}, §42, ``Quid Est Veritas?'' (1918), pp.~95--96.} Mencken recruits Paul's question and ``ways past finding out'' for a comic-polemical creed: the most ecstatic religion, he says, trusts without professing to know. Apostolic doxology is reversed into a case for agnosticism; that conclusion is Mencken's, not Paul's.

\item \textbf{``Go ye therefore, and teach all nations'' --- David Walker, \emph{Appeal}, art.~3 (1829), p.~42.} Walker confronts white American Christians with the missionary mandate and asks what they actually taught Africans: warfare, kidnapping, and slave trading rather than baptismal discipleship. The Great Commission becomes an antislavery indictment in which the supposed teachers stand accused by their own commission.
\end{itemize}
}

\Needspace{20\baselineskip}
\section*{The Propers: Interpretive Possibilities}

\noindent\emph{These are bounded AI-assisted editorial discovery proposals, attributed to none of the cited authorities. Each joins precisely named appointed elements by a stated mechanism; none claims historical selection intent, universal novelty, or doctrinal authority.}

\Needspace{9\baselineskip}
\noindent\textbf{Mercy becomes a missionary way of knowing.\enspace\properrefs{Int., Ep., Gosp., Off., Comm.}}
Tobit's command to bless God for shown mercy frames the Mass at entrance, offering, and Communion; between those repetitions, Romans reaches wonder before unsearchable wisdom and the Gospel sends disciples to all nations. The mechanism is a conversion of received action into epistemic humility and then public witness: mercy is known by grateful confession, not possessed as an explanation. Read element by element, one misses how the recurring Tobit text carries Paul's wonder outward into mission. The fruit is a rule for theology and evangelization---speak truly and confidently, but as a beneficiary rather than a proprietor of mystery. Experience does not measure revealed doctrine, thanksgiving is not a missionary technique, and no compositional link among Tobit, Romans, and Matthew is asserted.

\Needspace{9\baselineskip}
\noindent\textbf{The abyss is known asymmetrically.\enspace\properrefs{Ep., Grad., Gosp.}}
Romans calls God's judgments unsearchable; Daniel blesses the Lord who beholds the abysses while seated above the cherubim; Matthew gives the Church the singular Name of Father, Son, and Holy Ghost. The connecting mechanism is an asymmetry of vision and disclosure: creatures cannot climb from an abyss to comprehension, yet the one who sees every depth can stoop to name himself truthfully. Separate readings leave either apophatic wonder or a baptismal formula; their conjunction discovers why mystery can be confessed without being mastered. The intellectual and spiritual fruit is disciplined confidence in inquiry, prayer, and catechesis. Daniel's abysses are not Paul's judgments, grammar alone is not a complete Trinitarian treatise, and this contrast is no hidden scriptural cross-reference.

\Needspace{9\baselineskip}
\noindent\textbf{The Name moves from font to offering.\enspace\properrefs{Gosp., Sec., Postcomm.}}
The Gospel joins baptism to the singular Name of the three Persons; the Secret invokes God's holy Name over the victim and asks that the offerers become an eternal gift; the Postcommunion asks that confession advance salvation in body and soul. The mechanism is ritual transposition: a Name first received as belonging becomes invocation, self-offering, and embodied perseverance. Read separately, initiation, oblation, and post-Communion growth remain discrete; together they disclose baptismal identity as a vocation to be offered, not merely possessed. The pastoral fruit is a concrete baptismal examen of body, work, and allegiance. The Secret does not repeat Matthew's threefold grammar, the offerers are not identical without remainder with Christ the victim, and human self-gift neither causes nor completes revelation.

\Needspace{9\baselineskip}
\noindent\textbf{``From, through, and unto'' becomes a liturgy of return.\enspace\properrefs{Ep., Sec., Comm.}}
Romans confesses all things as from God, through God, and unto God; the Secret places victim and offerers before him; Communion turns received mercy into blessing before all the living. The mechanism is doxological return: creaturely existence received from God is carried through sacrificial mediation into grateful self-presentation and public praise. Element-by-element reading can leave Paul's prepositions as metaphysics and the prayers as devotion; their conjunction makes dependence ethically and liturgically answerable. Its fruit is a demanding stewardship of possessions, gifts, and speech: nothing is self-originating or finally self-directed. Romans does not distribute the three prepositions mechanically among the Persons, does not by itself command this exact oblation, and the creature's return remains grace-enabled through Christ rather than a closed natural cycle.

\Needspace{9\baselineskip}
\noindent\textbf{Doctrine acquires duration in the body.\enspace\properrefs{Coll., Gosp., Postcomm.}}
The Collect asks that true confession be guarded by steadfast faith; the Gospel commands baptism, teaching, and observance; the Postcommunion asks that sacramental reception and confession advance salvation in body and soul. The mechanism is a chain of temporal verbs---guard, teach, observe, advance---by which doctrine becomes a durable form of life. Separate readings can isolate orthodoxy, mission, and sacramental fruit; together they expose catechesis that does not reach practice and bodies as unfinished. The fruit is a test for formation: can a confessed truth become habit, corporate discipline, and mercy in bodily life? Moral performance neither completes revelation nor earns grace, and ``health of body and soul'' promises no exemption from illness, suffering, or gradual growth.
